\documentclass{homework}

\name{Student Name} % Replace (Student Name) with your name.
\id{}
\course{Seminars on Quantum Computing}
\hwnum{3}
\usepackage{tikz}

%\hwname{(Name)}          % Uncomment and replace (Name) with the type of the
                          % homework (e.g, Assignment, Problem Set, etc.) if you
                          % don't want the document to be labeled as "Homework."
%\problemname{(Name)}     % Uncomment and replace (Name) with the desired label
                          % for problems created with the problem environment.
%\solutionname{(Name)}    % Uncomment and replace (Name) with the desired label
                          % for solutions created with the solution environment.

% Load any other packages you need here.

\begin{document}

In this problem set, the following notations are used

\begin{equation*}
  \sigma_{x} =
  \begin{pmatrix}
    0 & 1\\ 1 & 0
  \end{pmatrix},\quad
  \sigma_{y} =
  \begin{pmatrix}
    0 & -i\\
    i & 0
  \end{pmatrix},\quad
  \sigma_{z} =
  \begin{pmatrix}
    1 & 0\\
    0 & -1
  \end{pmatrix},\quad
  \ket{+} = \frac{\ket{0} + \ket{1}}{\sqrt{2}}.
\end{equation*}

\begin{problem}
  Let $\rho$ be the density matrix of a qubit and let
  $U_{a, b} = \sigma_{z}^{a}\sigma_{x}^{b}$.
  The one-time pad $\mathrm{OTP}_{a,b}(\rho)$ with keys $a, b\in \{0,1\}$ is
  defined to be
  \begin{equation*}
    \mathrm{OTP}_{a, b}(\rho) = U_{a, b} \rho U_{a, b}^{\dagger}.
  \end{equation*}

  \begin{parts}
    \part\label{1.a} Prove that for any qubit state $\rho$,
    \begin{equation*}
      \E_{a, b} \mathrm{OTP}_{a, b}(\rho) = \frac{I}{2}.
    \end{equation*}

    \part\label{1.b} It is possible to do interesting computation on the state
    $\rho$ protected by the one-time pad without decoding.
    Show that if we apply an $X$ gate on $\mathrm{OTP}_{a, b}(\rho)$, we get
    $\mathrm{OTP}_{a, b}(X \rho X)$.

    \part\label{1.c} Show that if we apply $H$ on $\mathrm{OTP}_{a, b}(\rho)$ we obtain the
    one-time pad state $\mathrm{OTP}_{b, a}(H\rho H)$.
    Notice that the value of $a, b$ are swapped in this case!
  \end{parts}
\end{problem}

\begin{solution}

  \ref{1.a} % chktex 2
  Write the solution here.

  \ref{1.b} % chktex 2

  \ref{1.c} % chktex 2

\end{solution}

\begin{problem}
  Assume that state is either $\ket{\psi_{0}}$ or $\ket{\psi_{1}}$ with equal
  probability and $\abs{\langle \psi_{0} | \psi_{1} \rangle} = s > 0$.
  Design a positive-operator valued measure (POVM) $\{M_{0}, M_{1}, M_{*}\}$
  with three outcomes $0$, ${1}$, and $*$ to distinguish the state such that
  \begin{enumerate}
    \item if the outcome is $0$, the state must be $\ket{\psi_{0}}$,
    \item if the outcome is $1$, the state must be $\ket{\psi_{1}}$,
    \item the probability of having $*$ as the outcome is the smallest possible.
          Note that in this case, the measurement does not tell us whether the
          state is $\ket{\psi_{0}}$ or $\ket{\psi_{1}}$, and it's allowed to
          answer ``I don't know''.
          So the goal here is to minimize the probability of outputting an
          inconclusive answer.
  \end{enumerate}
\end{problem}

\begin{solution}

\end{solution}

\begin{problem}
  Let $\rho$ be a two-qubit state.
  Prove that the process of (1) applying a CNOT gate using the first qubit as
  the control and the second qubit as the target and (2) measuring the second
  qubit in $Z$ basis is the same as measuring the state $\rho$ directly using
  observable $Z\otimes Z$.
\end{problem}

\begin{solution}

\end{solution}

% \begin{problem}
%   Recall that for qubit-qubit entangled states, the PPT condition is a
%   sufficient and necessary condition for separability.
%   Use the PPT condition to compute the value of the smallest positive real
%   number $p$ such that
%   \begin{equation*}
%     \rho_{p} = p \frac{I}{4} + (1-p)\ket{\beta_{00}}\bra{\beta_{00}}
%   \end{equation*}
%   is a separable state.
% \end{problem}

% \begin{solution}

% \end{solution}

\begin{problem}
  In this problem, we consider a game between a quantum computer and a classical
  referee.
  Let $f : {\{0,1\}}^n \to {\{0,1\}}^{n-1}$ be a $2$-to-$1$ function, meaning
  that every $y \in {\{0,1\}}^{n-1}$ has exactly two distinct preimages.
  Suppose there is an efficient (with a polynomial number of gates) quantum
  circuit to compute $f$, but no efficient circuit that can produce from given
  $x \in {\{0,1\}}^{n}$ an $x' \ne x$ such that $f(x') = f(x)$.
  Show that, like in the Simon's algorithm, the quantum computer can efficiently
  generate a uniformly random $y$ and an associated $n$-qubit state
  $\ket{\phi_{y}}$.
  The quantum computer sends the string $y$ to the classical referee and keeps
  $\ket{\phi_{y}}$.
  Next the referee flips a coin and asks the quantum computer to perform one of
  the following two tasks.
  \begin{enumerate}
    \item (Task 0).
          When asked, efficiently generate an $x$ such that $f(x) = y$; and
    \item (Task 1).
          When asked, efficiently sample uniformly from the set
          $\{a \in {\{0,1\}}^{n} \mid a \cdot (x \oplus x') = 0 \mod 2\}$, where
          $x, x'$ are the two preimages of $y$ under $f$.
  \end{enumerate}
  Show that the quantum computer can always succeed.

  \noindent
  \emph{(Note: This game is used in a proof of quantumness protocol as classical
  computers cannot always succeed.)}
\end{problem}

\begin{solution}

\end{solution}

\begin{problem}
  Let \(U\) be an \(n\)-qubit unitary operator.
  Consider the input state 
  \begin{equation*}
    \rho_{\mathrm{in}} = \ket{0}\bra{0} \otimes \frac{I}{2^n}.
  \end{equation*}
  Apply the following operations to the state:
  \begin{enumerate}
    \item Apply a Hadamard gate \(H\) to the first qubit.
    \item Apply a controlled-\(U\) gate, where the first qubit is the control
          qubit.
    \item Apply another Hadamard gate \(H\) to the first qubit.
    \item Measure the first qubit in the computational basis.
  \end{enumerate}
  Define success as obtaining the measurement outcome \(0\) on the first qubit.
  Show that the success probability is
  \begin{equation*}
    P(0) = \frac{1}{2}
    \left(
      1+ \operatorname{Re} \left[
        \frac{\operatorname{Tr}(U)}{2^n} \right]
    \right).
  \end{equation*}
\end{problem}

\begin{solution}

\end{solution}

\end{document}
